\section{The Id Was Already There}
\label{sec:ch18-the-id-was-already-there}

\index{trace id!sent to every subgraph|(}%
Section~\ref{sec:ch18-the-page-that-got-slower} found one trace id in five logs
and called it the router's. That undersells it. The router did not only write
the id down; it handed a copy to every service it asked, on every fetch, and
the services took it.

\index{traceparent@\code{traceparent}!what the router sends}%
Anything that records request headers, placed between the router and the
services by pointing \code{graph.yaml}'s routing urls at it, shows what
arrives. Here is one run of the same page, one line per subgraph request, each
reduced to the name of the service it went to and the one header that
matters:

\begin{minted}[fontsize=\footnotesize]{text}
sessions  traceparent: 00-515457460e4ab1a45a57b5b745491f0d-856a85a0e4c3f7c4-01
speakers  traceparent: 00-515457460e4ab1a45a57b5b745491f0d-d32570b5eacb6961-01
ratings   traceparent: 00-515457460e4ab1a45a57b5b745491f0d-3b798df673db39f6-01
\end{minted}

\index{W3C Trace Context}%
\code{traceparent} is a header defined by the W3C Trace Context recommendation,
and its value is four fields separated by hyphens: a version, a trace id of
sixteen bytes written as thirty-two hex characters, a parent id of eight bytes
written as sixteen, and a byte of flags~\autocite{w3ctracecontext}. The trace
id is the request's and does not change as the request travels. The parent id
is the identity of the call being made, and the specification requires every
participant to replace it on each request it sends onward, which is why the
three lines above share their second field and differ in their third.

\index{Cosmo Router!propagates trace context with no configuration}%
None of that was configured. The \code{config.yaml} this book has printed since
chapter~\ref{ch:entities-authorization} has no \code{telemetry} block at all,
there is no graph token, and the router's startup log lists
\enquote{Cosmo Cloud Tracing} among the features it has disabled for want of
one. The id is minted and propagated anyway. WunderGraph's reference for the
router says that it enables the Trace-Context specification by
default~\autocite{cosmotelemetry}, and that is the whole mechanism: the router
does not need somewhere to send traces in order to give each request an id and
pass it on.

\subsection{What the Services Did With It}

\index{ASP.NET Core!adopts an incoming trace id}%
\index{Activity@\code{Activity}}%
Each service received that header and did not ignore it. ASP.NET Core reads an
incoming \code{traceparent} and makes the request's
\code{System.Diagnostics.Activity}, the runtime's object for \enquote{the
operation in progress}, a continuation of it. Microsoft documents that the
ASP.NET web server decodes and encodes these ids on HTTP messages with no
code from the application~\autocite{msdistributedtracing}, and it holds for
the services in this book exactly as chapter~\ref{ch:ownership} left them:
inside each one, \code{Activity.Current} carries the router's trace id, and
its parent id is the one the router put on that service's fetch.
Section~\ref{sec:ch18-writing-it-down} makes each service print the trace id.

\index{logging!scopes carry the trace id and are not printed}%
The same id is already attached to everything those services log. .NET's
logging adds a \emph{scope}, a set of values stamped onto every log entry
written while it is active, carrying \code{TraceId}, \code{ParentId} and \code{SpanId}, the last being the
service's own id for the work it is doing, which it would pass on as the
parent id of any call it made. Those three are on by
default~\autocite{msaspnetlogging}.
The console's default formatter does not print scopes. And
\code{StatementLog} bypasses the logger entirely, writing its marker with
\code{Console.WriteLine}, so even a formatter that printed scopes would not
have put the id on the three lines
section~\ref{sec:ch18-the-page-that-got-slower} could not attribute.

\subsection{Where the Id Stops}

\index{trace id!where it stops}%
Figure~\ref{fig:ch18-one-id-three-fetches} is that request drawn with what
each component does with the id. The router mints it, writes it down and sends
it three times. Each service receives it, carries it through the request, and
writes it nowhere. And the response goes back to the client without it.

\begin{figure}[htbp]
  \centering
  % One request, one trace id, three fetches that never see it recorded
% anywhere but the router.
%
% Every other figure in this book is a timeline read left to right, one
% lane per scenario (figures/tikz/ch01-one-field.tex, ch09-two-fetches.tex,
% ch10-one-fetch-two-costs.tex). This chapter's argument is about five
% parties in the same request rather than about stages within one lane, so
% this figure turns the axis: five lifelines run top to bottom instead of
% one axis running left to right, and the sequence reads down the page. The
% idiom itself is reused unchanged: square corners, one 0.4pt stroke
% throughout, the same -{Straight Barb} arrowhead as every axis and span in
% this book, no fill anywhere, and the same solid/dashed split ch09 and
% ch10 already use for a request versus its return. Lifelines are dashed
% for the same reason ch09's seam boundary is dashed: a guide the eye
% follows without mistaking it for a message.
%
% The router mints the trace id itself; nothing upstream sent one. That
% mint is not a message between two lifelines, so it is drawn as a tick on
% the router's own lifeline with a solid-outline label beside it, not an
% arrow. Solid outline is deliberate and is the figure's one write/no-write
% distinction: the router's box below ("logs T") is the same solid outline
% as this mint note, and every other party's box ("holds T, writes
% nothing", "never sees T") is dashed. Solid outline means this party wrote
% the id down; dashed means it did not.
%
% Sessions answers before the router fans out to Speakers and Ratings.
% Both fetches leave from the router's own lifeline within a few mm of each
% other rather than from one literal shared point: two coincident straight
% lines cannot both be drawn without one running on top of the other for
% the 72mm they would otherwise share, so the arrow to Ratings starts 6mm
% below the arrow to Speakers. Both labels sit left-aligned just to the
% right of the router lifeline, stacked at that same offset, so neither
% overlaps the other and neither reads as attached to a lifeline further
% along the row.
%
% No response carries the trace id back. The two return arrows from
% Speakers and Ratings, and the final response to the client, are dashed
% like every return in this book and carry no traceparent label, because
% none of them has one to show.
%
% Every node whose content varies (the five lifeline heads, the five
% write/hold boxes) carries a leading \vphantom{Tg} so its height and depth
% come from that fixed reference glyph rather than from its own letters.
% Without it, a word with a descender (speakers, ratings, writes) measures
% taller than one without (client, router, sessions, logs), and anchoring
% on south or north then puts its box at a visibly different level. This
% is the same fix under every node here, not a special case for one row.
\begin{tikzpicture}[
  x=1mm, y=1mm,
  every node/.append style={font=\sffamily\scriptsize},
  head/.style={draw, line width=0.4pt, align=center, inner sep=2pt,
    minimum width=18mm, anchor=south},
  lifeline/.style={draw, line width=0.4pt, dashed},
  request/.style={draw, line width=0.4pt,
    -{Straight Barb[length=1.4mm, width=1.6mm]}},
  reply/.style={draw, line width=0.4pt, dashed,
    -{Straight Barb[length=1.4mm, width=1.6mm]}},
  msglabel/.style={anchor=south, align=center, inner sep=1pt},
  msglabelleft/.style={anchor=south west, align=left, inner sep=1pt},
  tick/.style={draw, line width=0.4pt},
  mint/.style={draw, line width=0.4pt, align=center, inner sep=1.5pt,
    anchor=west},
  writesbox/.style={draw, line width=0.4pt, align=center, inner sep=2pt,
    text width=20mm, anchor=north},
  holdsbox/.style={draw, line width=0.4pt, dashed, align=center,
    inner sep=2pt, text width=20mm, anchor=north},
]

% ------------------------------------------------------------- lifelines
\node[head, anchor=south west] at (0,8)   {\vphantom{Tg}client};
\node[head]                    at (36,8)  {\vphantom{Tg}router};
\node[head]                    at (72,8)  {\vphantom{Tg}sessions};
\node[head]                    at (108,8) {\vphantom{Tg}speakers};
\node[head, anchor=south east] at (144,8) {\vphantom{Tg}ratings};

\foreach \x in {0,36,72,108,144}{
  \draw[lifeline] (\x,3) -- (\x,-90);
}

% --------------------------------------------------------- 1. the query
\draw[request] (0,-14) -- (36,-14);
\node[msglabel, text width=30mm] at (18,-12)
  {POST /graphql\\{\itshape\tiny no trace id}};

% --------------------------------------------------- 2. the router mints
\draw[tick] (34,-24) -- (38,-24);
\node[mint] at (39,-24) {trace id T};

% ------------------------------------------------------- 3. to sessions
\draw[request] (36,-34) -- (72,-34);
\node[msglabel] at (54,-32) {traceparent 00-T-p1-01};

\draw[reply] (72,-44) -- (36,-44);

% ------------------------- 4. sessions answered: fan out, a few mm apart
\draw[request] (36,-54) -- (108,-54);
\node[msglabelleft] at (37,-52) {00-T-p2-01};

\draw[request] (36,-60) -- (144,-60);
\node[msglabelleft] at (37,-58) {00-T-p3-01};

\draw[reply] (108,-68) -- (36,-68);
\draw[reply] (144,-74) -- (36,-74);

% --------------------------------------------------------- 5. the reply
\draw[reply] (36,-84) -- (0,-84);
\node[msglabel, text width=30mm] at (18,-82) {200, no trace id};

% ----------------------------------------------------- what gets written
\node[holdsbox, anchor=north west] at (0,-94)
  {\vphantom{Tg}never sees T};
\node[writesbox]                   at (36,-94)
  {\vphantom{Tg}logs T};
\node[holdsbox]                    at (72,-94)
  {\vphantom{Tg}holds T, writes nothing};
\node[holdsbox]                    at (108,-94)
  {\vphantom{Tg}holds T, writes nothing};
\node[holdsbox, anchor=north east] at (144,-94)
  {\vphantom{Tg}holds T, writes nothing};

\end{tikzpicture}

  \caption{One page, one trace id. The router attaches the id to each of its
  three fetches with a fresh parent id, and every service holds it for the
  length of the request. Only the router writes it anywhere, and the client,
  who is the one person able to say which request was slow, is never told
  it.}
  \label{fig:ch18-one-id-three-fetches}
\end{figure}

\index{trace id!a client's own trace is continued}%
A client that already has a trace of its own, a browser or a mobile app carrying its own
instrumentation, can send a \code{traceparent} to the router, and the router
continues it: its log line and every subgraph fetch then carry the client's
trace id rather than one the router made up. So the join is available from
the very first process in the chain, if the first process asks for it.

\index{config version!which composition answered}%
\code{config\_version}, on the same log line, is the field
chapter~\ref{ch:enter-the-router} could not make mean anything. It is
thirty-two zeroes in this book because the execution
config is composed on a laptop. In a deployment that fetches its config from a
control plane it is the version of the composition that served the request,
which answers the question every slow-page investigation eventually reaches:
was this before or after the last publish.

\index{Netflix!routing a question to its owner from a trace}%
This is not a problem peculiar to a four-service conference graph. Tejas
Shikhare, describing Netflix's federated graph, gives the reason the response
alone cannot carry the blame: \enquote{In GraphQL, almost every response is a
200 with custom errors in the error block}. His team's answer to what he calls
the \enquote{who do I ask about\ldots} routing problem was to link types and
fields to their owning teams' support channels, so that finding one \enquote{is
now as simple as clicking a link from a trace}~\autocite{shikhare2020part2}. A
link from a trace needs a trace, and a trace starts from an id that every
service wrote down.
\index{trace id!sent to every subgraph|)}%
